% Options for packages loaded elsewhere
\PassOptionsToPackage{unicode}{hyperref}
\PassOptionsToPackage{hyphens}{url}
\PassOptionsToPackage{space}{xeCJK}
\documentclass[
  chinese,
  11pt,
]{ctexart}
\usepackage{xcolor}
\usepackage[margin=2.2cm]{geometry}
\usepackage{amsmath,amssymb}
\setcounter{secnumdepth}{5}
\usepackage{iftex}
\ifPDFTeX
  \usepackage[T1]{fontenc}
  \usepackage[utf8]{inputenc}
  \usepackage{textcomp} % provide euro and other symbols
\else % if luatex or xetex
  \usepackage{unicode-math} % this also loads fontspec
  \defaultfontfeatures{Scale=MatchLowercase}
  \defaultfontfeatures[\rmfamily]{Ligatures=TeX,Scale=1}
\fi
\usepackage{lmodern}
\ifPDFTeX\else
  % xetex/luatex font selection
  \ifXeTeX
    \usepackage{xeCJK}
    \setCJKmainfont[]{Noto Serif SC}
    \setCJKsansfont[]{Noto Sans SC}
    \setCJKmonofont[]{DejaVu Sans Mono}
  \fi
  \ifLuaTeX
    \usepackage[]{luatexja-fontspec}
    \setmainjfont[]{Noto Serif SC}
  \fi
\fi
% Use upquote if available, for straight quotes in verbatim environments
\IfFileExists{upquote.sty}{\usepackage{upquote}}{}
\IfFileExists{microtype.sty}{% use microtype if available
  \usepackage[]{microtype}
  \UseMicrotypeSet[protrusion]{basicmath} % disable protrusion for tt fonts
}{}
\makeatletter
\@ifundefined{KOMAClassName}{% if non-KOMA class
  \IfFileExists{parskip.sty}{%
    \usepackage{parskip}
  }{% else
    \setlength{\parindent}{0pt}
    \setlength{\parskip}{6pt plus 2pt minus 1pt}}
}{% if KOMA class
  \KOMAoptions{parskip=half}}
\makeatother
\usepackage{longtable,booktabs,array}
\newcounter{none} % for unnumbered tables
\usepackage{calc} % for calculating minipage widths
% Correct order of tables after \paragraph or \subparagraph
\usepackage{etoolbox}
\makeatletter
\patchcmd\longtable{\par}{\if@noskipsec\mbox{}\fi\par}{}{}
\makeatother
% Allow footnotes in longtable head/foot
\IfFileExists{footnotehyper.sty}{\usepackage{footnotehyper}}{\usepackage{footnote}}
\makesavenoteenv{longtable}
\ifLuaTeX
\usepackage[bidi=basic,shorthands=off,provide=*]{babel}
\else
\usepackage[bidi=default,shorthands=off,provide=*]{babel}
\fi
\ifLuaTeX
  \usepackage{selnolig} % disable illegal ligatures
\fi
\setlength{\emergencystretch}{3em} % prevent overfull lines
\providecommand{\tightlist}{%
  \setlength{\itemsep}{0pt}\setlength{\parskip}{0pt}}
\usepackage{bookmark}
\IfFileExists{xurl.sty}{\usepackage{xurl}}{} % add URL line breaks if available
\urlstyle{same}
\hypersetup{
  pdftitle={课题04-0.1B端到端},
  pdflang={zh-CN},
  hidelinks,
  pdfcreator={LaTeX via pandoc}}

\title{课题04-0.1B端到端}
\author{}
\date{}

\begin{document}
\maketitle

{
\setcounter{tocdepth}{2}
\tableofcontents
}
\section{\texorpdfstring{课题04 开题简报 · 0.1B 端到端预训练（\textbf{首次上机}）}{课题04 开题简报 · 0.1B 端到端预训练（首次上机）}}\label{ux8bfeux989804-ux5f00ux9898ux7b80ux62a5-0.1b-ux7aefux5230ux7aefux9884ux8badux7ec3ux9996ux6b21ux4e0aux673a}

\begin{quote}
模块：\textbf{M4/M5} ｜ 周次：\textbf{W5--W6（10-12→10-25）} ｜ 算力：\textbf{T2 Kaggle 2×T4，计划 12--21 GPU·小时} 唐杰作业对应：\textbf{① 端到端训一个 0.1B} ｜ 判分来源：\textbf{自建}（loss/bpb 达标 + 采样可读性）+ \texttt{nanochat/core\_eval.py} 的 CORE 手写：复用课题01/02 的产物（tokenizer + 模型 + 优化器），本课题新增\textbf{数据管线与训练编排的判断}
\end{quote}

\begin{center}\rule{0.5\linewidth}{0.5pt}\end{center}

\subsection{一、研究背景}\label{ux4e00ux7814ux7a76ux80ccux666f}

``0.1B'' 是唐杰作业里的数字，也是\textbf{单卡/双卡免费算力下的上限}。参考点：nanochat 的 d12 约等于 GPT-1 量级 （官方讨论区原话：``d12 是我最喜欢的模型，它是 GPT-1 的大小，约 6 分钟训完''------那是在 8×H100 上）。 \textbf{本课题的核心不是刷指标，而是把''训练一次完整模型''的每个环节亲自走通并记账。}

\subsection{二、研究问题（一句话）}\label{ux4e8cux7814ux7a76ux95eeux9898ux4e00ux53e5ux8bdd}

\begin{quote}
\textbf{在 2×T4、20 GPU·小时的预算内，我能不能从零训出一个''能生成连贯英文、并在 CORE 上有非零分''的 0.1B 模型？实测成本与估算差多少？}
\end{quote}

\subsection{三、攻关线索}\label{ux4e09ux653bux5173ux7ebfux7d22}

\begin{enumerate}
\def\labelenumi{\arabic{enumi}.}
\tightlist
\item
  \textbf{数据}：TinyStories（全量）+ owt-sample 采样。\textbf{先算一遍}：词表大小 × 语料字节 → token 数 → 需要的步数
\item
  \textbf{规模选择}：nanochat 用单一旋钮 \texttt{-\/-depth}；在 T4 上 16GB 显存能放多深？（用 M0 的显存公式先估）
\item
  \textbf{精度}：\textbf{T4 是 SM75，没有 bf16} → 必须 fp16 + GradScaler（或用 nanochat 的 \texttt{NANOCHAT\_DTYPE=float16}）
\item
  \textbf{断点续训}：Kaggle 会话 9--12h 会断 → checkpoint 必须写 \texttt{/kaggle/working}（20GB 持久）
\item
  \textbf{超时自停}：脚本内置 wall-clock 上限，防止忘记关会话
\item
  \textbf{指标}：训练看 bpb（词表无关），最终跑 CORE（22 数据集合成，\texttt{core\_eval.py} 单文件）
\end{enumerate}

\subsection{四、里程碑}\label{ux56dbux91ccux7a0bux7891}

\begin{itemize}
\tightlist
\item[$\square$]
  M1 \textbf{实验卡获批}：写清假设/命令/预算上限/停止条件/成功判据（见《算力与成本预算.md》§五）
\item[$\square$]
  M2 \textbf{会话预检}：\texttt{nvidia-smi} 确认 2×T4（\textbf{若是 P100 立即重开会话}）
\item[$\square$]
  M3 \textbf{冒烟}：\texttt{-\/-depth=4}、100 步，确认能跑、能存、能续训（≤10 分钟）
\item[$\square$]
  M4 首次正式训练：双 T4 + DDP，跑到预算上限或 loss 平台
\item[$\square$]
  M5 续训一次（验证 checkpoint 真的可用 ------ 这一条最容易翻车）
\item[$\square$]
  M6 评测：bpb + CORE + \textbf{采样 20 条样本}（人工判断''是否连贯''）
\item[$\square$]
  M7 \textbf{成本账单}：实测耗时 vs 估算（M0 的 6ND + T4 折算），逐项解释偏差
\item[$\square$]
  M8 一页报告 + 知识卡（含''我最意外的三个现象''）
\end{itemize}

\subsection{五、交付物}\label{ux4e94ux4ea4ux4ed8ux7269}

{\def\LTcaptype{none} % do not increment counter
\begin{longtable}[]{@{}
  >{\raggedright\arraybackslash}p{(\linewidth - 4\tabcolsep) * \real{0.3333}}
  >{\raggedright\arraybackslash}p{(\linewidth - 4\tabcolsep) * \real{0.3333}}
  >{\raggedright\arraybackslash}p{(\linewidth - 4\tabcolsep) * \real{0.3333}}@{}}
\toprule\noalign{}
\begin{minipage}[b]{\linewidth}\raggedright
交付物
\end{minipage} & \begin{minipage}[b]{\linewidth}\raggedright
位置
\end{minipage} & \begin{minipage}[b]{\linewidth}\raggedright
验收
\end{minipage} \\
\midrule\noalign{}
\endhead
\bottomrule\noalign{}
\endlastfoot
实验卡 & \texttt{证据/实验卡-NN.md} & 经用户确认 \\
run.yaml + 日志 & \texttt{证据/} & 含 commit、配置、随机种子、耗时 \\
曲线 & \texttt{证据/} & loss/bpb vs FLOPs 与 vs 时间 两张 \\
采样样本 & \texttt{证据/样本.md} & ≥20 条，标注''通顺/不通顺'' \\
\textbf{成本账单} & \texttt{报告.md} & 实测 vs 估算偏差 ≤ 2× 或给出解释 \\
checkpoint & Kaggle 持久盘 & 不进程（体积原因）；只在 \texttt{证据/} 留哈希与路径 \\
\end{longtable}
}

\subsection{六、讲课锚点}\label{ux516dux8bb2ux8bfeux951aux70b9}

\begin{itemize}
\tightlist
\item
  \textbf{讲义}：\texttt{M4-训练.md}（W3 展开）
\item
  \textbf{对标}：\texttt{nanochat/README.md} 与 \texttt{runs/speedrun.sh}（读它怎么组织全流程）· CS336 lecture 05/13--14
\item
  \textbf{搭桥}：课题03 的 scaling 结论在这里第一次被''真跑''检验
\end{itemize}

\subsection{七、代码级提问示例}\label{ux4e03ux4ee3ux7801ux7ea7ux63d0ux95eeux793aux4f8b}

\begin{enumerate}
\def\labelenumi{\arabic{enumi}.}
\tightlist
\item
  双 T4 跑 d12，为何\textbf{不是} 2 倍加速？先预测扩展效率（50\%？80\%？），再测。
\item
  T4 上 fp16 与 fp32，谁更快？谁更稳？GradScaler 在解决什么？
\item
  你的 loss 在第 2000 步开始上升。按 M0 §五 的顺序，你的前三个动作是什么？
\item
  若把语料从 TinyStories 换成中文语料（词表不变），loss 会怎么变？\textbf{bpb 呢？}
\end{enumerate}

\subsection{八、风险与提醒}\label{ux516bux98ceux9669ux4e0eux63d0ux9192}

\begin{itemize}
\tightlist
\item
  🚨 \textbf{P100 陷阱}：本课题如果拿到 P100，训练慢 3--5 倍且后续课题05 直接不可用 → \textbf{宁可等 T4 时段}。
\item
  🚨 \textbf{配额}：本课题占全学期约 1/5 的配额，实验卡必须写实；\textbf{冒烟不通过不许开正式训练}。
\item
  ⚠️ 0.1B 在 T4 上大概率需要''降 batch + 梯度累积''，这会改变有效 batch size → \textbf{要写进 run.yaml}。
\item
  ⚠️ Kaggle 会话到期会杀进程：\textbf{每 500 步存一次}是硬要求。
\end{itemize}

\end{document}
